% ---- score2 standard preamble (from: python3 score2_tools.py preamble) ----
\documentclass[10pt]{extarticle}
\usepackage[a4paper,margin=1.5cm]{geometry}
\usepackage[T1]{fontenc}
\usepackage{amsmath,amssymb}
\usepackage{graphicx}
\usepackage[export]{adjustbox}
\usepackage{multicol}
\usepackage{enumitem}
\usepackage{booktabs,array,tabularx}
\setlength{\columnsep}{0.7cm}
\setlength{\parindent}{0pt}
\setlength{\parskip}{3pt}
\pagestyle{empty}
\makeatletter
% figure registry (figdefs.tex is generated by the extract step)
\newcommand{\qfigdef}[4]{\@namedef{qf@file@#1}{#2}\@namedef{qf@w@#1}{#3}\@namedef{qf@wide@#1}{#4}}
\newcommand{\qf@check}[1]{\@ifundefined{qf@file@#1}{\PackageError{score2}{Unknown figure id #1}{}}{}}
\newcommand{\qf@img}[4][]{\includegraphics[width=#2\dimexpr\@nameuse{qf@w@#4}pt\relax,max width=#3,max height=0.45\textheight,#1]{\@nameuse{qf@file@#4}}}
% Full-width questions: finalize wraps every question that contains a wide figure in \qfullbegin...\qfullend,
% so the WHOLE question (stem, figures, options) is set across the page in reading order, then the two columns
% resume. The question's head - its number, stem and first figure (\qfullsplit marks the end of the head) -
% is kept on one page, so a number is never left at a page foot with its figure overleaf; the rest may break
% across pages normally, so no big gaps are left. Without a \qfullsplit, the whole question is kept together
% if it is under 45% of a page. Never write these yourself.
\newif\ifqf@full
\newif\ifqf@boxed
\newsavebox{\qf@box}
\newcommand{\qf@open}{\global\qf@fulltrue\par\global\qf@boxedtrue
  \setbox\qf@box\vbox\bgroup\hsize\textwidth\linewidth\textwidth}
\newcommand{\qf@emit}[1]{\par\egroup\global\qf@boxedfalse
  \ifdim\dimexpr\ht\qf@box+\dp\qf@box\relax<#1\textheight\noindent\usebox{\qf@box}\par\else\unvbox\qf@box\fi}
\newcommand{\qfullsplit}{\ifqf@boxed\qf@emit{0.85}\fi}
\newcommand{\qf@close}{\ifqf@boxed\qf@emit{0.45}\fi\global\qf@fullfalse}
\newcommand{\qfullbegin}{\end{multicols}\qf@open}
\newcommand{\qfullend}{\qf@close\begin{multicols}{2}}
\newcommand{\qfullstart}{\qf@open}
\newcommand{\qfulldone}{\qf@close}
% \qfig[scale]{id}: block figure, centred, directly where it is written (size set by the tool).
% Fallback if a wide figure is not inside a full-width question: break out of the columns in place - no floats.
\newcommand{\qfig}[2][1]{\qf@check{#2}%
  \ifqf@full\qf@blk{#1}{#2}\else
  \ifnum\@nameuse{qf@wide@#2}=1 \expandafter\@firstoftwo\else\expandafter\@secondoftwo\fi
  {\end{multicols}\begin{center}\qf@img{#1}{\textwidth}{#2}\end{center}\begin{multicols}{2}}%
  {\qf@blk{#1}{#2}}\fi}
\newcommand{\qf@blk}[2]{\par\smallskip{\centering\qf@img{#1}{\linewidth}{#2}\par}\smallskip}
% \qfigi[scale]{id}: inline figure (inside option lists / tables / minipages); never breaks columns.
\newcommand{\qfigi}[2][1]{\qf@check{#2}\qf@img[valign=c]{#1}{\linewidth}{#2}}
\makeatother
% structure macros - lint reads these to check completeness and figure placement
\newcommand{\papersec}[1]{\par\bigskip{\centering\Large\bfseries PAPER-#1\par}\medskip}
\newcommand{\partsec}[2]{\par\medskip\noindent{\large\bfseries PART-#1: #2}\par\smallskip}
\newcommand{\sect}[1]{\par\smallskip\noindent\textbf{#1}\par}
\newcommand{\qn}[1]{\par\smallskip\noindent\textbf{#1.}~}
\newcommand{\paraq}[2]{\par\smallskip\noindent\textbf{Paragraph for Questions #1 and #2}\par}
\newenvironment{qoptions}{\begin{itemize}[leftmargin=1.8em,itemsep=1pt,topsep=2pt,label={}]}{\end{itemize}}
\newcommand{\opt}[2]{\item[(#1)] #2}
% figure registry (generated by score2_tools.py finalize)
\qfigdef{p1_phy_s1_q01_1}{Score-I - Test 19 (20 Nov 2022)-img/p1_phy_s1_q01_1.png}{119.6}{0}
\qfigdef{p1_phy_s1_q02_1}{Score-I - Test 19 (20 Nov 2022)-img/p1_phy_s1_q02_1.png}{108.2}{0}
\qfigdef{p1_phy_s1_q03_1}{Score-I - Test 19 (20 Nov 2022)-img/p1_phy_s1_q03_1.png}{108.3}{0}
\qfigdef{p1_phy_s1_q04_1}{Score-I - Test 19 (20 Nov 2022)-img/p1_phy_s1_q04_1.png}{125.4}{0}
\qfigdef{p1_phy_s1_q05_1}{Score-I - Test 19 (20 Nov 2022)-img/p1_phy_s1_q05_1.png}{161.1}{0}
\qfigdef{p1_phy_s1_q06_1}{Score-I - Test 19 (20 Nov 2022)-img/p1_phy_s1_q06_1.png}{108.3}{0}
\qfigdef{p1_phy_s2_q01_1}{Score-I - Test 19 (20 Nov 2022)-img/p1_phy_s2_q01_1.png}{153.1}{0}
\qfigdef{p1_phy_s2_q02_1}{Score-I - Test 19 (20 Nov 2022)-img/p1_phy_s2_q02_1.png}{42.6}{0}
\qfigdef{p1_phy_s2_q03_1}{Score-I - Test 19 (20 Nov 2022)-img/p1_phy_s2_q03_1.png}{92.5}{0}
\qfigdef{p1_phy_s2_q05_1}{Score-I - Test 19 (20 Nov 2022)-img/p1_phy_s2_q05_1.png}{92.5}{0}
\qfigdef{p1_phy_s2_q06_1}{Score-I - Test 19 (20 Nov 2022)-img/p1_phy_s2_q06_1.png}{159.0}{0}
\qfigdef{p1_phy_s3_q01_1}{Score-I - Test 19 (20 Nov 2022)-img/p1_phy_s3_q01_1.png}{171.7}{0}
\qfigdef{p1_phy_s3_q02_1}{Score-I - Test 19 (20 Nov 2022)-img/p1_phy_s3_q02_1.png}{126.1}{0}
\qfigdef{p1_phy_s3_q03_1}{Score-I - Test 19 (20 Nov 2022)-img/p1_phy_s3_q03_1.png}{78.7}{0}
\qfigdef{p1_phy_s3_q04_1}{Score-I - Test 19 (20 Nov 2022)-img/p1_phy_s3_q04_1.png}{108.2}{0}
\qfigdef{p1_phy_s3_q05_1}{Score-I - Test 19 (20 Nov 2022)-img/p1_phy_s3_q05_1.png}{152.0}{0}
\qfigdef{p1_phy_s3_q06_1}{Score-I - Test 19 (20 Nov 2022)-img/p1_phy_s3_q06_1.png}{192.5}{0}
\qfigdef{p1_chem_s1_q04_1}{Score-I - Test 19 (20 Nov 2022)-img/p1_chem_s1_q04_1.png}{232.7}{0}
\qfigdef{p1_chem_s2_q01_2}{Score-I - Test 19 (20 Nov 2022)-img/p1_chem_s2_q01_2.png}{313.9}{1}
\qfigdef{p1_chem_s2_q02_1}{Score-I - Test 19 (20 Nov 2022)-img/p1_chem_s2_q02_1.png}{129.4}{0}
\qfigdef{p1_chem_s2_q06_1}{Score-I - Test 19 (20 Nov 2022)-img/p1_chem_s2_q06_1.png}{273.1}{1}
\qfigdef{p1_chem_s3_q02_1}{Score-I - Test 19 (20 Nov 2022)-img/p1_chem_s3_q02_1.png}{215.1}{0}
\qfigdef{p1_chem_s3_q05_1}{Score-I - Test 19 (20 Nov 2022)-img/p1_chem_s3_q05_1.png}{319.9}{1}
\qfigdef{p1_chem_s3_q06_1}{Score-I - Test 19 (20 Nov 2022)-img/p1_chem_s3_q06_1.png}{107.8}{0}
% ---- body skeleton ----
% \begin{document}
% \begin{center}{\Large\bfseries JEE (Main+Advanced) Score-II --- Test NN (DD Mon YYYY)}\end{center}
% \begin{multicols}{2}
%   \papersec{1}   (only for two-paper tests)
%   \partsec{1}{PHYSICS} \sect{Section-I (i): Single Correct} \qn{1} ... \qfig{p1_phy_q01_1} ...
% \end{multicols}
% \end{document}

\begin{document}
\begin{center}{\Large\bfseries JEE (Main+Advanced) Score-I --- Test 19 (20 Nov 2022)}\end{center}
\raggedcolumns
\setlength{\tabcolsep}{3pt}
\begin{multicols}{2}
\papersec{1}
\partsec{1}{PHYSICS}
\sect{SECTION-I}
\qn{1} Spring is in its natural length. Collision of the block with the ring fixed with the horizontal surface and the wall is perfectly elastic. If at $t=0$, a velocity $v_0=0.5$ m/s is given to the block and spring constant of spring is 100 N/m, mass of block is 1 kg, then :-
\qfig{p1_phy_s1_q01_1}
\begin{qoptions}
\opt{A}{0.3 m/s is velocity of block just after collision with ring.}
\opt{B}{0.4 m/s is velocity of block just after collision with ring.}
\opt{C}{$\pi/10$ sec is time period of the motion.}
\opt{D}{0.4 m/s is velocity of block just before collision with wall.}
\end{qoptions}

\qn{2} A long metal strip of thickness $t=1$ mm and width $d=1$ cm carries current along its length ($L$). The current in the metal strip is following towards the left. It is placed in a uniform transverse magnetic field $B=2$ T (as shown in figure). The potential difference recorded across the width `$d$' of the plate is found to be $3\times10^{-6}$ volt. Select the correct statement.
\qfig{p1_phy_s1_q02_1}
\begin{qoptions}
\opt{A}{Electric field resulted across the width `$d$' is $1\times10^{-3}$ V/m.}
\opt{B}{Electric field resulted across the width `$d$' is $3\times10^{-4}$ V/m.}
\opt{C}{Magnitude of electric force on the electron due to the electric field that is resulted due to the motion of the electron inside metal strip is $4.8\times10^{-22}$ N.}
\opt{D}{The drift velocity of the electron is $1.5\times10^{-4}$ m/s.}
\end{qoptions}

\qn{3} A conducting rod of mass `$m$' is placed on frictionless horizontal rails and a force of ``$mg$'' starts acting on the conducting rod at $t=0$ toward right perpendicular to rod. One end of the parallel conducting rod is connected by the inductor of inductance `$L$' and other end of the conducting rod is connected by the capacitor as shown in the diagram. A uniform vertical magnetic field exist in the region perpendicular to the plane of the ring. (take $CB^2\ell^2=$ mass of rod)
\qfig{p1_phy_s1_q03_1}
\begin{qoptions}
\opt{A}{The maximum horizontal displacement of the rod is $\dfrac{4mgL}{B^2\ell^2}$.}
\opt{B}{The maximum horizontal displacement of the rod is $\dfrac{MgL}{B^2\ell^2}$.}
\opt{C}{The angular frequency of the oscillation of the rod is $\dfrac{B\ell}{\sqrt{2mL}}$.}
\opt{D}{None}
\end{qoptions}

\qn{4} A circular region of radius $R=\sqrt3$ m which has a uniform magnetic $B=0.2$ Tesla directed into the plane of the figure. A particle having mass $m=2$ gm, speed $v=0.3$ m/s and charge $q=1$ mC is projected along the radius of the circular region as shown in diagram. Select the correct statement(s).
\qfig{p1_phy_s1_q04_1}
\begin{qoptions}
\opt{A}{Deviation produced in the charge particle is $\dfrac{\pi}{3}$.}
\opt{B}{Deviation produced in the charge particle is $\dfrac{\pi}{2}$.}
\opt{C}{Average acceleration of the charge particle is $9\pi$ cm/s$^2$.}
\opt{D}{Distance travelled by the particle in the magnetic field is $\pi$ m.}
\end{qoptions}

\qn{5} A metallic disc of radius `$a$' is mounted on a rigid axle, is rotating with constant angular velocity $\omega$ inside a long solenoid of inductance $L$ whose two ends are connected to the rotating disc of by two brush contact as shown. Resistivity of the wire of solenoid is $\rho$ and resistance $R$ is connected between solenoid and metallic disc as shown in figure. Then :-
\qfig{p1_phy_s1_q05_1}
\begin{qoptions}
\opt{A}{The induced e.m.f. between terminals P and Q is $\dfrac{\mu_0 Ni\omega a^2}{2\ell}$ when current in the coil is $i$.}
\opt{B}{Current in the coil increases or decrease with time depending upon the value of $\omega$.}
\opt{C}{$\omega_{\min}$ for increasing the current in coil is $\dfrac{\ell R}{2\mu_0 N\rho a^2}$.}
\opt{D}{$\omega_{\min}$ for increasing the current in the coil is $\dfrac{2\ell R}{\mu_0 N a^2}$.}
\end{qoptions}

\qn{6} A thin walled non-magnetic spherical shell with wall thickness $d$ and radius $R$ ($R\gg d$) is placed in a uniform magnetic field $\vec B=B_0(\hat j)$, which slowly changes as $\dfrac{dB}{dt}=\alpha$, where $\alpha$ is a positive constant. Electrical resistivity of the material of the shell is `$\rho$'. Then select the correct option.
\qfig{p1_phy_s1_q06_1}
\begin{qoptions}
\opt{A}{Magnetic moment induced in the shell is in negative `$y$' direction.}
\opt{B}{Magnetic moment induced is directly proportional to $d$.}
\opt{C}{Magnetic moment induced is inversely proportional to $d$.}
\opt{D}{Magnetic moment is independent of the $d$.}
\end{qoptions}

\sect{SECTION-II}
\qn{1} A pulse source of microwave produces burst of 20 GHz radiation with each burst having a duration 1 ns. A parabolic reflector ($R=6$ cm) is used to produce a parallel beam of radiation as indicated in the diagram. The average output power of the source during each pulse is 25 kW. If microwave travel in air, maximum value of electric field $E$ is $x\times10^{4}$ V/m. Then find the value of $x$.
\qfig{p1_phy_s2_q01_1}

\qn{2} Four identical massless springs of the spring constant 70 N/m, two massless pulley and a block of mass of 6.6 kg are connected as shown and system is kept in vertical plane. Initially the system is in equilibrium. Now the block is slightly displaced from equilibrium position of equilibrium position along a vertical line and then released. Find the time period (in sec).
\qfig{p1_phy_s2_q02_1}

\qn{3} An infinity long current carrying wire is placed at the origin along the $z$-axis. Current in the wire is 2 amp. A straight wire $AB$ of length 2 m is placed from $A(2,0)$ to $B\left(3,\sqrt3\right)$ as shown in diagram. Then the value of $\displaystyle\int_A^B \vec B\cdot d\vec\ell$ is $x\times10^{-8}$ (in S.I. Unit). Find the value of $x$.
\qfig{p1_phy_s2_q03_1}

\qn{4} Suppose Bundi is a small town near Kota with a demand of 1200 kW electric power at 220 Volt is located 30 km from KTPS (Kota Thermal Power Station) generating power at 440 Volt. The resistance of the two vires line carrying power is 0.2 $\Omega$/km. Bundi gets power from the line through a 4000 V -- 220 Volt step down transformer at substation in Bundi. Assuming no power loss in the transformer. Power supplied by the KTPS to Bundi substation is :- (in kW)

\qn{5} The 8-kg body is moved 0.2 m to the right of the equilibrium position and released from rest at time $t=0$. Determine angular frequency of the damped oscillation. (Given that) The viscous damping coefficient $c$ is 20 N$\cdot$s/m, and the spring stiffness $k$ is 32 N/m.
\qfig{p1_phy_s2_q05_1}

\qn{6} A coil is wound on a hollow insulating cylinder, which contains in it a laminated iron core. The dependence of inductance $L$ with the displacement $x$ of the iron core is shown in the figure. In the initial state $x=0$ (the core is fully inserted into the coil), the current in the coil is 1.0 A. Find current in the coil immediately after the core is quickly taken out of the coil. (Answer in Amperes)
\qfig{p1_phy_s2_q06_1}

\sect{SECTION-III}
\qn{1} The circuit shown in figure is driven by an AC voltage $16\sin\left(\dfrac{5t}{\sqrt3}\right)$ volts. The capacitor (area of each plate of capacitor $=10^{-4}$ m$^2$) is suddenly removed at $t=\dfrac{7\pi}{10\sqrt3}$ seconds and then the geometrically, identical dielectric slabs of dielectric constant of $K_1=2$ and $K_2=6$ respectively. Find surface charge density at the interface in kC/m$^2$. If your answer is $\alpha$ then find the value of $\alpha/150$.
\qfig{p1_phy_s3_q01_1}

\qn{2} A hollow conducting cone has mass $10\pi$ gm and height 20 cm and radius 5 cm. A current $I$ flowing uniformly along its curved surface area. This cone is kept on a very rough surface and a uniform magnetic field $\vec B=4\hat i$ tesla exist in space. The maximum value of $I$ (in Amp) so that the cone does not topple about point P is :-
\qfig{p1_phy_s3_q02_1}

\qn{3} In the AC circuit shown in the diagram. $L=80$ mH and $C=2.5$ mF and $R=6\,\Omega$. The rms voltage of AC source is 50 volt of frequency 100 rad/sec. The rms current (in amp) through the capacitor is $I_0$. Find the value of $I_0/2$.
\qfig{p1_phy_s3_q03_1}

\qn{4} When magnetic flux through a coil is changed the variation of induced current in the coil with time is as shown in the graph. If resistance of coil is $2\,\Omega$, then the magnitude of total change in flux (in weber) of coil will be :
\qfig{p1_phy_s3_q04_1}

\qn{5} A homogeneous field of magnetic induction $B$ is perpendicular to a track of gauge $l$ which is inclined at an angle $\alpha=30^\circ$ to the horizontal. A frictionless conducting rod of mass $m$ straddles the two rails of the track as shown in the figure. If maximum speed of the conductor is $=\dfrac{mgR}{nB^2l^2}$. Then find $n$
\qfig{p1_phy_s3_q05_1}

\qn{6} Two long conducting bars P and Q are fixed parallel to each other in a horizontal plane some distance apart. On these bars, two parallel jumpers $J_1$ and $J_2$ each of resistance $R$ are arranged at separation $x_0$ as shown in the figure. The jumpers can slide on the bars without friction. At some point of time, a uniform vertical magnetic field of induction $B$ is switched on. If the final separation between the jumpers is $\dfrac{x_0}{N}$, find $N$. The force between the jumpers due to their currents is negligible as compared to the force of interaction between current in the jumper and the magnetic field. (Self inductance of system is neglected)
\qfig{p1_phy_s3_q06_1}

\partsec{2}{CHEMISTRY}
\sect{SECTION-I}
\qn{1} Which of the following statement(s) is/are CORRECT ?
\begin{qoptions}
\opt{A}{Aspirin is an example of non-narcotic analgesics.}
\opt{B}{Soap is formed by boiling triglyceryl stearate with caustic soda}
\opt{C}{Saccharin and alitame have S atom in their structure.}
\opt{D}{In HDPE preparation, Zeigler natta catalyest is used up}
\end{qoptions}

\qn{2} Which of the following statement(s) is/are CORRECT ?
\begin{qoptions}
\opt{A}{Vitamin C is water soluble}
\opt{B}{Cytosine, Thymine and Uracil are bases derived from purine.}
\opt{C}{A nucleic acid is a polynucleotide by condensing sugar, phosphate and nitrogeneous base where the polymeric link is a phosphodiester.}
\opt{D}{Adenine and thymine are paired by two hydrogen bonds whereas Guanine and Cytosine are paired by three hydrogen bonds.}
\end{qoptions}

\qn{3} Which of the following statement(s) is/are NOT correct ?
\begin{qoptions}
\opt{A}{Duma's method is not applicable to compounds containing nitrogen in nitro and azo groups.}
\opt{B}{Sulphur is determined in an organic compound by heating it in a carius tube with Na$_2$O$_2$ and precipitating all sulphur as BaSO$_4$.}
\opt{C}{For determining phosphorus, the organic compound is heated with conc. HNO3 which oxidizes all phosphorus to H$_3$PO$_4$. It is precipitated as ammonium phosphomolybdate by adding NH$_3$ and ammonium molybdate.}
\opt{D}{Oxygen is estimated by decomposing the organic compound by heating it in a stream of N$_2$ gas. The gaseous mixture is passed over red hot coke which converts all O$_2$ to CO. The CO produced is treated with warm I$_2$O$_5$ which produces CO$_2$ and I$_2$. The CO$_2$ produced gives estimation of oxygen.}
\end{qoptions}

\qn{4} Consider the following sequence of reactions :
\qfig{p1_chem_s1_q04_1}
A gives positive DNP test. B gives positive NaHCO$_3$ test.\\
Which of the following is/are matched correctly ?
\begin{qoptions}
\opt{A}{A on reaction with NaOD in D$_2$O shows an increase of 3 amu in molecular mass.}
\opt{B}{The solution of B in benzene gives colligative properties with van't Hoff factor greater than 1.}
\opt{C}{Z on reaction with benzaldehyde in presence of CH$_3$COONa followed by acidification gives a product with molecular mass 148 amu.}
\opt{D}{The pattern of hydrolysis of C and D in NaOH solution can be different}
\end{qoptions}

\qn{5} Which of the following reaction(s) will result in formation of Me$_3$C--NH--CH$_3$
\begin{qoptions}
\opt{A}{Me$_3$C--NH--CHO $\xrightarrow[\text{(ii) H}_2\text{O}]{\text{(i) LAH, dry ether}}$}
\opt{B}{Me$_3$C--NH$_2$ $\xrightarrow{\text{CHCl}_3+\text{KOH}}$ $\xrightarrow[\text{(ii) H}_2\text{O}]{\text{(i) LAH, dry ether}}$}
\opt{C}{Me$_3$C--CO--CH$_3$ + CH$_3$NH$_2$ $\longrightarrow$}
\opt{D}{Me$_2$C=CH$_2$ $\xrightarrow{\text{CH}_3\text{NH}_2}$}
\end{qoptions}

\qn{6} Consider the following compounds :\\
(I) CCl$_3$CHO \quad (II) CH$_3$CHO \quad (III) CH$_3$COCH$_3$ \quad (IV) Ph--CO--CH$_3$\\
Select the CORRECT statement(s) :
\begin{qoptions}
\opt{A}{Extent of hydration : I $>$ II $>$ III $>$ IV}
\opt{B}{II \& III on reaction with CaOCl$_2$ solution followed by heating of the Ca containing moeity give back the respective starting compound.}
\opt{C}{All compounds show positive iodoform test.}
\opt{D}{III is more polar than II.}
\end{qoptions}

\qfullbegin
\sect{SECTION-II}
\qn{1} In the following roadmap, 78 g of benzene is used at the start. Percentage yields, wherever applicable, are indicated in brackets. At the end of all the steps, what mass of the product is obtained ?\\
(Atomic mass data : C = 12, H = 1, N = 14, O = 16, Cl = 35.5, Br = 80)
\qfig{p1_chem_s2_q01_2}
\qfullend

\qn{2} \qfigi{p1_chem_s2_q02_1}\\
Given the ratio of the molar mass of R in g/mol to the degree of unsaturation present in the compound

\qn{3} A 2 L, 0.1 M solution of sodium acetate is electrolyzed using Pt electrodes. If a current of 9.65 mA is passed for 2500 seconds, and the current efficiency is 80\%, then the mass of hydrocarbon in mg produced at the anode will be-\\
(1 faraday = 96500 C)

\qn{4} The number of chiral center present in maltose

\qn{5} Among the twenty naturally occuring amino acids, if\\
$a$ = number of amino acids whose pH at isoelectric point is $<4$.\\
$b$ = number of amino acids whose pH at isoelectric point is $>7$.\\
$c$ = number of amino acids which have an aromatic ring in their structure.\\
Find the value of expression $a\times\dfrac{b}{c}$.

\qfullbegin
\qn{6} In the following reaction sequence, the amount of D (in g) formed from 10 moles of acetophenone is\underline{\hspace{2em}}. (Atomic weight in g mol$^{-1}$: H = 1, C = 12, N = 14, O = 16, Br = 80. The yield (\%) corresponding to the product in each step is given in the parenthesis)
\qfig{p1_chem_s2_q06_1}
\qfullend

\sect{SECTION-III}
\qn{1} Consider the heating reactions of dicarboxylic acids, HOOC--(CH$_2$)$_x$--COOH where $x$ is the number of methylene groups between the COOH groups varies from 1 to 5.\\
$p$ = number of acids in which one of the products is CO$_2$\\
$q$ = number of acids in which one of the products in H$_2$O\\
$r$ = number of acids in which products contain both CO$_2$ and H$_2$O\\
Calculate ($p+q-r$)

\qn{2} \qfigi{p1_chem_s3_q02_1}\\
If the molar mass of S is $x$ g/mol, report $x/17$.

\qn{3} How many of the following drugs are bactericidal in nature?\\
Erythromycin, tetracycline, penicillin, ofloxacin, chloramphenicol.

\qn{4} How many of the following are addition polymers ?\\
Polyethene, Teflon, PVC, Nylon, Bakelite, Melamine-formaldehyde resin, PMMA, Polystyrene, Natural rubber, Urea-formaldehyde resin.

\qfullbegin
\qn{5} In how many of the following compounds is the reactivity greater than that of benzene towards nitration via electrophilic aromatic substictution mechanism.
\qfig{p1_chem_s3_q05_1}
\qfullend

\qn{6} In a developed chromatogram of a flower extract shown, find the retardation factor ($x$) of the segment. Report $5x$.
\qfig{p1_chem_s3_q06_1}

\partsec{3}{MATHEMATICS}
\sect{SECTION-I}
\qn{1} The coefficient of $x^6$ in the expansion of $(1+x^2)^3\left(1-\sqrt[3]{x}\right)^7\left(1+\sqrt[3]{x}\right)^4$ is :
\begin{qoptions}
\opt{A}{49}
\opt{B}{$-43$}
\opt{C}{43}
\opt{D}{0}
\end{qoptions}

\qn{2} 15 boys and 5 girls are to be seated in a row randomly. The probability that between any two girls atleast 2 boys sit is
\begin{qoptions}
\opt{A}{$\dfrac{6!\,7!}{20!\,19!}$}
\opt{B}{$\dfrac{12!\,15!}{6!\,20!}$}
\opt{C}{$\dfrac{15!\,12!}{7!\,20!}$}
\opt{D}{$\dfrac{16!\,18!}{20!\,7!}$}
\end{qoptions}

\qn{3} The value of $\dfrac{1}{2^{15}}\displaystyle\sum_{r=0}^{15}(-1)^r\,5^{15-r}\;{}^{30}C_r\;{}^{30-r}C_{15-r}$ is
\begin{qoptions}
\opt{A}{number of ways of selection of 15 persons from 30 married couples such that no husband is selected with his wife}
\opt{B}{number of ways of selection of 15 persons from 15 married couples}
\opt{C}{number of ways of writing answers to 15 true-false type questions selected out of 30 questions}
\opt{D}{number of ways of writing answers to 15 multiple choice correct (one or more choice correct) questions selected out of 30 questions. (each question has 4 choices)}
\end{qoptions}

\qn{4} Consider a set `$S$' of 15-digit numbers ($d_i\in\{0,1,2,\ldots,9\}$)\\
$S=\{d_1d_2d_3\ldots d_{14}d_{15}\mid d_i\le d_{i+1}\}$ for all $i\in\{1,2,\ldots,14\}$\\
Let $n(S)$ denotes the number of elements in $S$, then
\begin{qoptions}
\opt{A}{exponent of 3 in $n(S)$ is 3}
\opt{B}{exponent of 3 in $n(S)$ is 1}
\opt{C}{unit digit in $(n(S))^4$ is 0}
\opt{D}{unit digit in $(n(S))^4$ is 6}
\end{qoptions}

\qn{5} Which of the following statement(s) is/are INCORRECT ?
\begin{qoptions}
\opt{A}{$\log_{(\sqrt6-\sqrt5)}\left(\sqrt8-\sqrt7\right)>1$}
\opt{B}{$\log_{(\sqrt6-\sqrt5)}\left(\sqrt8-\sqrt7\right)<1$}
\opt{C}{$\log_7\left(\sqrt[3]{2}+\sqrt[3]{5}\right)\left(\sqrt[3]{4}-\sqrt[3]{10}+\sqrt[3]{25}\right)$ is an integer value}
\opt{D}{number of rational terms in the expansion of $\left(7^{\frac17}+2^{\frac15}\right)^{105}$ is 3}
\end{qoptions}

\qn{6} If `$x$' satisfies the equation $\displaystyle\sum_{r=1}^{2023}\sin(rx)\cos(2024-r)x=\dfrac{2023}{2}$ then
\begin{qoptions}
\opt{A}{$\sin6\theta=\sin(2024x)$ is true for 6 real values of $\theta$ in $[0,2\pi]$}
\opt{B}{$\sin6\theta=\sin(2024x)$ is true for 12 real values of $\theta$ in $[0,2\pi]$}
\opt{C}{$\sin(2024x)>\cos(2024x)$}
\opt{D}{$\sin(2024x)<\cos(2024x)$}
\end{qoptions}

\sect{SECTION-II}
\qn{1} Let `$a$' be the number of ways in which 3 consecutive persons can be selected from $n$ persons sitting around a round table and `$b$' be number of ways of selecting one male and one female from $n+3$ males and $n+3$ females, then the value of $36\left(\displaystyle\sum_{a=4}^{2023}\dfrac{a+3}{(a+1)!\cdot b}+\dfrac{1}{2026!}\right)$ is

\qn{2} Number of value(s) of $x$ in $[0,2\pi]$ satisfying the equation\\
$\left(\log_{(\sqrt2-1)}\tan x\right)$\\
$\sqrt{\log_{(\sqrt2-1)}\sqrt{(3-2\sqrt2)\left(\sqrt2-1\right)}+\log_{\tan x}\left(\sqrt2-1\right)}=-1$ is

\qn{3} Number of real solution(s) of the equation $\dfrac{\sin^22\theta}{2}+6\sin\theta\cdot\cos\theta(\sin\theta+\cos\theta)=10$ in $[0,4\pi]$ is

\qn{4} If $\left(\log_pq+\log_qp\right)^2=2029$, $p>q>1$, then the number of values of $\theta\in(0,6\pi)$ satisfying the equation $\dfrac{1}{\log_{pq}q}-\dfrac{1}{\log_{pq}p}=\operatorname{cosec}\theta$ is

\qn{5} Two friends `$A$' and `$B$' plays a game by picking one ball alternately without replacement from a bag containing 7 white and 4 red balls in the order ABABAB\ldots.. The one who gets a first red ball wins the game. The probability that A got two attemps to draw balls given that B wins the game is $\dfrac mn$ (in lowest terms, $m,n\in I$) then $n-m$ is

\qn{6} A random variable $X$ has the probability distribution\par\smallskip
\begin{center}
\renewcommand{\arraystretch}{1.1}
\begin{tabular}{|c|c|c|c|c|c|c|c|c|c|}
\hline
\textbf{X} & 1 & 2 & 3 & 4 & 5 & 6 & 7 & 8 & 9\\
\hline
\textbf{P(X)} & 0.11 & 0.2 & 0.14 & 0.26 & 0.13 & 0.01 & 0.03 & 0.04 & 0.08\\
\hline
\end{tabular}
\end{center}
for the events $A=\{X$ is an odd number$\}$,\\
\hspace*{2em}$B=\{X$ is the exponent of 13 in $100!\}$,\\
\hspace*{2em}$C=\{X$ is unit digit of $7^n$, $n\in N\}$, the probability $P((A\cap B)\cup C)$ is

\sect{SECTION-III}
\qn{1} Let $X$ be a binomial variate with mean and variance 30 and 5 respectively and $Y$ be a binomial variate with mean and variance 9 and $\dfrac{27}{4}$ respectively. Let $P_1(X=r)$ and $P_2(Y=r)$ denotes the probabilities of obtaining `$r$' success for random variables $X$ and $Y$ respectively. Let $a,b$ be two positive integers (different from one) satisfying $a^{a+b}=b^{12}$ and $b^{a+b}=a^3$ and if $\dfrac{P_1(X=a)}{P_2(Y=b)}=\dfrac{11\cdot17\cdot2^{35}\cdot p^4}{q^{70}}$ then $p-q$ is

\qn{2} Consider the following sets\\
$A=\{\theta:\sec\theta+\operatorname{cosec}\theta=\sin\theta+\cos\theta,\ \theta\in[0,2\pi]\}$,\\
$B=\left\{\theta:\cos\theta=\sqrt{1-\sin\theta},\ \theta\in[0,3\pi]\right\}$,\\
$C=\Bigl\{\theta:2^{\sin^2\theta}>\dfrac{\cos\theta}{2\cos\theta+1},$\\
\hspace*{1em}$\theta\in\mathbb R\Bigr\}$. The number of elements in the set $(A\cup B)\cap C$ is

\qn{3} Number of real solution(s) of equation $\dfrac13\log_2(\cos^35x+\cos^37x)=1+\log_2\cos6x+\log_2\cos x$ in $\left[0,\dfrac\pi2\right]$ is

\qn{4} Consider the observations\\
\hspace*{1em}$x_i=\dfrac52+p,\ i=1$ to $20$,\\
\hspace*{1em}$x_i=\dfrac72,\ i=21$ to $40$,\\
\hspace*{1em}$x_i=\dfrac52-p,\ i=41$ to $60$,\\
if the variance of these 60 observations is $\dfrac49$ then the value of $6p^2$ is

\qn{5} Let $f(m)=\displaystyle\sum_{r=1}^{m}\dfrac{\left(m^r+m^{r-1}\right)\left(1-\frac1r\right)}{{}^mC_r}$ then the unit digit of the value of $f(2023)$ is

\qn{6} Let $N$ be the number of ways in which a kid can arrange 5 balls, one each of red, green, black, yellow and pink in a row such that, yellow is not adjacent to black and pink. Green is not adjacent to black and red. The value of $\dfrac27N$ is
\end{multicols}
\end{document}
